% Folder 142, batch 03 — archivists' pages 41–60.
% Pass: Opus 5.5 (claude-opus-5-5), 2026-10-02 — first pass, unchecked against the pages by a human.
%
% Page 41 is the tail of a run (author's p. 23) begun before this batch;
% page 42 is a blank sheet carrying only a pencil tally and is skipped.
% From page 43 a new run starts, paginated by the author 1–11 (p. 43 = 1,
% p. 60 = 11), and it runs on past the end of the batch.
% Notation fixed for the batch: the author underlines the letters that name
% two-element sets (orientations, components); they are set in bold here.
\documentclass[11pt,a4paper]{article}
\input{../preamble/grothendieck}

\folder{142}
\batch{3}
\pages{41}{60}
\foldertitle{Action de Galois sur Teichmüller : notes manuscrites (s.d.).}
\dating{[à partir de 1978]}
\watermark{Édition de démonstration}

\begin{document}

\page{41}
\note{p. 23 de l'auteur. La page achève une suite de calculs commencée avant ce lot : les formules (76) à (79) auxquelles elle renvoie n'y figurent pas. Le reste du feuillet est blanc.}

l'égalité du $\beta'$ défini par cette équation \add{les formules équivalentes,} avec
celui défini par (76), ou (76$'$) ou (77), ou \struck{le fait que $\beta'$ (vu}
l'équation (78), écrite en choisissant pour $\beta'$ défini par l'une des formules
(76) à (77), ou par la formule (79) — ce qui donne la formulation la plus simple.

\section{La sphère à trois points marqués}
\note{Nouvelle suite, paginée 1 à 11 par l'auteur sur les pages 43 à 60 de ce lot ; elle se poursuit au-delà. Les lettres qu'il souligne pour désigner des ensembles à deux éléments ($\boldsymbol{\varpi}$, $\boldsymbol{\omega}$, $\boldsymbol{\varepsilon}$) sont rendues en gras.}

\page{43}
\note{p. 1 de l'auteur.}
\marginal{\uncertain{Attention} : ne pas confondre la notation \ill{} avec celles des \ill{} $\alpha_i, \beta_i, \beta'_i, \gamma_i, \gamma'_i$ \ill{} (1)…}

Considérons une 2-sphère conforme $\Sigma$ avec une partie $J$ à trois éléments,
\[
  J \subset \Sigma, \qquad J = (R_i)_{i \in J}
\]
de sorte que le groupe $\mathfrak{S}_J$ opère sur $\Sigma$ par automorphismes
\struck{orientés} directs, \struck{et aussi} \uncertain{étant} entendu que le groupe des
automorphismes \struck{$\mathfrak{S}_J \times \{1,\tau\}$ le groupe des} \add{conformes} de la
situation $(\Sigma, J)$ est
\[
  \Gamma_{\Sigma,J} = \Gamma \simeq \mathfrak{S}_J \times \{1, \tau\}
\]
où $\tau$ est l'unique \struck{anti}automorphisme $\neq 1$ \struck{\ill{}} qui fixe les
$R_i$ \struck{\ill{}} ($\tau$ = conjugaison complexe).
\marginal{NB L'antipodisme ne fait pas partie de $\Gamma$, car il ne stabilise pas $J$.}
\marginal{\ill{} dans la catégorie des $(\Sigma, J)$ \uncertain{équivalente} (2) \ill{} $(J, \boldsymbol{\varpi})$, $\boldsymbol{\varpi} \in (\mathrm{Ens}_2)$ \ill{}, et $\boldsymbol{\varpi} \in (\mathrm{Ens}_2)$ \ill{} $J$ plus bas.}
Il y a sur $\Sigma$ une unique métrique riemannienne homogène invariante
\add{à scalaire $\lambda \in \mathbf{R}^{*+}$ près} par $\mathfrak{S}_J$,
\struck{par laquelle}, qui nous permet d'identifier canoniquement $\Sigma$ à la sphère
unité dans un espace vectoriel réel euclidien $E_3$. On pose
\[
  (3)\quad
  \left\{
  \begin{array}{ll}
    \boldsymbol{\varpi} = \Omega(\Sigma) \simeq \Omega(E_3)
      & \text{ens. des deux orientations} \\
      & \text{de } \Sigma \text{, ou de } E \\
    \boldsymbol{\omega} = \Omega(J) \simeq \Omega(\Sigma_{\mathbf{R}}) & \\
    \boldsymbol{\varepsilon} = \pi_0(\Sigma \setminus \Sigma_{\mathbf{R}})
      = \pi_0(E \setminus V) = \pi_0\big((E/V)^{*}\big) &
  \end{array}
  \right.
\]
\marginal{On pourrait écrire $V(J, \boldsymbol{\omega})$, $E(J, \boldsymbol{\varpi})$.}
\marginal{(3 bis) On a $\boldsymbol{\omega} \wedge \boldsymbol{\varepsilon} \wedge \boldsymbol{\varpi} = 1$, i.e. \ill{} — pour la comm. des orientations \uncertain{voir plus loin} (4).}
où $V$ est le $\mathbf{R}$-sous-espace vectoriel de $E$ engendré par les $R_j$, qui est
de dimension 2, et où
\[
  \Sigma_{\mathbf{R}} = \Sigma \cap V
  \qquad \text{cercle euclidien du plan euclidien } V .
\]

\page{44}
On aura dans $V$
\[
  (4)\quad \sum_{j \in J} R_j = 0
\]
ce qui donne un isomorphisme
\[
  (5)\quad V \simeq V_J(\mathbf{R})
  = \operatorname{Ker}\big(\mathbf{R}^J \xrightarrow{\ \text{somme}\ } \mathbf{R}\big)
\]
et on a, vu la formule (3)
\[
  (6)\quad \Delta = V^{\perp} \simeq \mathbf{R}(\boldsymbol{\varepsilon})
\]
(isom. compatible avec la structure quadratique), $\Delta$ « axe polaire », orth. de $V$
dans $E$, donc
\[
  (7)\quad E \simeq V_J(\mathbf{R}) \oplus \mathbf{R}(\boldsymbol{\varepsilon})
  \qquad \text{isom. d'espaces euclidiens}
\]
où dans $V_J(\mathbf{R})$ on prend la seule forme quadratique $q$ pour laquelle les $R_j$
\struck{soient unitaires} satisfont
\struck{de carré un}
\[
  q(R_j) = 1
\]
(si \ill{} $q$ \ill{} prend une base formée de $R_i, R_j$, \ill{} $R_k$ de coordonnées
$(-1,-1)$, la forme quadratique en question est $x^2 - xy + y^2$).
\marginal{Il est clair, \ill{} $q$ \ill{} \uncertain{invariante} sous l'opération de $\mathfrak{S}_J$ sur $V(J)$.}
Dans le plan euclidien $V$, les $R_i$ forment les sommets d'un triangle équilatéral inscrit
dans le cercle unité. L'opération de $\mathfrak{S}_J$ sur $E$ se fait par son opération
\struck{\ill{}} tautologique sur $V(J)$, et l'opération par $\operatorname{sg}(g)$ sur
$\Delta$ :

\page{45}
\note{p. 2 de l'auteur.}
\[
  (8)\quad g_E = \underbrace{g_{V_J(\mathbf{R})}}_{\text{tautol.}}
  \oplus \operatorname{sg}(g)\, \mathrm{id}_{\Delta}
\]
L'opération de $\tau$ est donnée par
\[
  (9)\quad \tau_E = \mathrm{id}_{V_J(\mathbf{R})} \oplus (-\mathrm{id}_{\Delta})
\]
Je pose
\[
  (10)\quad \Sigma^{*} = \Sigma \setminus J = \Sigma - \{R_i \mid i \in I\}
\]
\note{$I$ pour $J$, tel quel sur la page.}

et je m'intéresse au groupoïde fondamental de $\Sigma^{*}$, comme exercice préliminaire
de syntaxe galoisienne, avant de regarder des situations \add{un peu} plus complexes
(le triangle équilatéral remplacé par d'autres configurations riches en symétries sur
la sphère, voire sur d'autres surfaces que la sphère, — celle-ci étant remplacée par
\uncertain{la fibre de Hopf}, en vue des applications au calcul des groupoïdes de
Teichmüller…).

Il est commode d'introduire une compactification de $\Sigma^{*}$, pour disposer aussi de
pts-base « proches » de chaque $R_i$.

\page{46}
\marginal{La construction de l'éclatement \uncertain{réel} et de la contraction \ill{} dans les variétés $C^\infty$ et \ill{} analytiques réelles.}
Pour ceci considérons l'éclaté réel de $\Sigma$ en $J$, soit $\Sigma'$, dans lequel
l'image inverse de $i \in J$ s'identifie à \struck{l'espace} la droite projective réelle
\[
  (11)\quad \Sigma'_{R_i} \simeq P(T_{\Sigma,R_i})
\]
($T_{\Sigma,R_i}$ espace tangent à $\Sigma$ en $R_i$, identifié à $T^{*}_{\Sigma,R_i}/\mathbf{R}^{*}$),
qui est un cercle plongé dans la surface compacte $\Sigma'$ (qu'on peut considérer comme
obtenue en enlevant de la sphère trois rondelles autour des $R_i$, et y recollant des
rubans de Möbius — on obtient une surface non orientable dont le \struck{genre} $\chi_!$
est égal à
\[
  \chi_!(\Sigma) - 3\chi_!(\text{rondelle ouverte})
  + 3\big(\underbrace{\chi_!(\text{ruban Möbius})}_{0}\big)
  = 2 - 3 = -1 = 1 - g ,
\]
donc le genre est $2$). Le découpage de cette surface suivant $\bigcup \Sigma'_{R_i}$ est une
surface à bord compacte $\widetilde{\Sigma}$, avec $\partial\widetilde{\Sigma}$ un
rev. double de $\coprod \Sigma'_{R_i}$, \struck{donc} soit en fait
\[
  (12)\quad \partial\widetilde{\Sigma} \simeq \coprod_{i \in I} TU_{\Sigma,R_i}
\]
(cercle tangent unitaire à $\Sigma$ en $R_i$, identifiable à
$T^{*}_{\Sigma,R_i}/\mathbf{R}^{*+}$), noté aussi $U_i$.
\marginal{$\widetilde{\Sigma}$ est une sphère à bord à trois trous. $\widetilde{\Sigma}^{\circ}$ \ill{} munie d'une structure conforme, mais celle-ci ne se prolonge pas à $\widetilde{\Sigma}$.}

\page{47}
\note{p. 3 de l'auteur.}
On a des équivalences de catégories
\[
  (13)\quad
  \operatorname{Rev\,et}(\widetilde{\Sigma})
  \xrightarrow[\text{équiv}]{\ \approx\ }
  \operatorname{Rev\,et}(\widetilde{\Sigma}^{\circ})
  \xrightarrow[\text{iso}]{\ \sim\ }
  \operatorname{Rev\,et}(\Sigma^{*})
\]
déduites de l'inclusion
\[
  \widetilde{\Sigma}^{\circ}
  \overset{\mathrm{def}}{=} \widetilde{\Sigma} \setminus \partial\widetilde{\Sigma}
  = \widetilde{\Sigma} \,|\, \Sigma^{*} \hookrightarrow \widetilde{\Sigma}
\]
et de l'isomorphisme
\[
  \widetilde{\Sigma}^{\circ} \xrightarrow{\ \sim\ } \Sigma^{*},
\]
d'où une équivalence entre les groupoïdes fondamentaux au sens large (cf. les objets
\uncertain{sont} les rev. étales 1-connexes)
\[
  (14)\quad \Pi_1(\widetilde{\Sigma}) \longrightarrow \Pi_1(\Sigma^{*})
\]
et en composant avec
\[
  \Pi_1(U_i) \longrightarrow \Pi_1(\widetilde{\Sigma})
\]
déduit de l'inclusion $U_i \hookrightarrow \widetilde{\Sigma}$, on trouve
\[
  (15)\quad \Pi_1\Big(\coprod U_i\Big) \longrightarrow \Pi_1(\Sigma^{*}) .
\]
Nous allons représenter graphiquement l'homomorphisme 14, en « dessinant »
$\widetilde{\Sigma}$ sur $\Sigma^{*}$, comme une région qui serait contenue dans
$\Sigma^{*}$ en lui enlevant \add{l'intérieur de} trois disques \add{fictifs} $\Delta_i$
« infiniment petits » (s'identifiant aux disques unités des $T_{\Sigma,R_i}$, ou encore
aux cônes sur les $U_i$), et on se représentera (14) comme provenant d'une telle
inclusion fictive.
\marginal{Plus précisément, au lieu de travailler \ill{} \ill{} trois trous \ill{} sur $\widetilde{\Sigma}$ \ill{} sphère \ill{} plus grande \ill{}.}

\page{48}
En fait, nous allons travailler avec des éléments de $\Pi_1(\widetilde{\Sigma})$ qui
proviennent, soit de pts de $\widetilde{\Sigma}$, donc de pts de $\Sigma^{*}$ ou d'un des
$U_i$, soit de parties de $\widetilde{\Sigma}$ qui sont 1-connexes pour la top. induite par
$\widetilde{\Sigma}$ — en partant \uncertain{peut-être} de parties \add{1-connexes} de
$\Sigma^{*}$, et en les faisant « déborder » sur $\widetilde{\Sigma}$ de façon naturelle.

Le groupe $\Gamma$ opère sur la situation \add{$C^\infty$} $(\Sigma, J)$, donc par
transport de structure sur la situation
\[
  \widetilde{\Sigma} \longrightarrow \Sigma, \qquad J \subset \Sigma
\]
donc sur (14), (15). Un aspect essentiel du travail à faire est de garder cette
\uncertain{opération} constamment présente à l'esprit, et d'introduire des notations et
des constructions qui soient invariantes, dans le sens requis par la situation, par les
opérations en question.

\page{49}
\note{p. 4 de l'auteur.}
\drawing{une sphère, pôles $P^{+}$ en haut et $P^{-}$ en bas, l'équateur $\Sigma_{\mathbf{R}}$ portant les trois trous $R_0$, $R_1$, $R_\infty$, chacun entouré de son petit cercle $U_i$ marqué des points $R_i^{\pm}$, $S_i^{\pm}$ et des arcs $a_i^{++}, a_i^{+-}, \dots$ ; sur l'équateur les points $Q_0$, $Q_1$, $Q_\infty$ et les arcs orientés $b_i^{\pm}$ ; des pôles partent les méridiens portant les arcs $c_i^{\pm}$ (vers les $Q_i$) et $d_i^{\pm}$ (vers les $S_i$), en traits pleins ou pointillés selon l'hémisphère.}

On introduit
\[
  (16)\quad
  \left\{
  \begin{array}{ll}
    Q_i \ (i \in J) & \text{antipodique de } R_i \ (= -R_i) \\
    R_i^{\omega} \in U_i & (i \in J,\ \omega \in \boldsymbol{\omega}) \\
    P_{\varepsilon} & (\varepsilon \in \boldsymbol{\varepsilon}
      = \pi_0(\Sigma \setminus \Sigma_{\mathbf{R}})) \\
    S_i^{\varepsilon} &
  \end{array}
  \right.
\]
où $P_\varepsilon$ est défini par
\[
  (17)\quad \{P_\varepsilon\} = \Delta \cap (\Sigma - \Sigma_{\mathbf{R}})_\varepsilon
\]
\struck{et $R_i^{\omega} \in U_{i\mathbf{R}}$ \ill{}} et $R_i^{\omega}$ est défini ainsi :
On a une bijection canonique
\[
  (18)\quad
  U_{i\mathbf{R}} \overset{\mathrm{def}}{=}
  \underbrace{U_i \cap T_{\Sigma_{\mathbf{R}},R_i}}_{TU_{\Sigma_{\mathbf{R}},R_i}}
  \simeq \boldsymbol{\omega}
\]
($\Sigma_{\mathbf{R}}$ variété de dim 1, l'une des orientations de celle-ci aux
\uncertain{points} $R_i$ (donc globale) est un \ill{} bijectif avec l'une des deux
directions orientées tangentes).
\[
  (19)\quad S_i^{\varepsilon} \in U_i \cap T_{\Sigma_i} = TU_{\Sigma_i}
\]
\struck{\ill{}}
où $\Sigma_i$ est le grand cercle déterminé par les $P_\varepsilon$ et $Q_i$ ; on a
$TU_{\Sigma_i} \simeq \boldsymbol{\varepsilon}$ (orientations de celle-ci, car les deux
directions tangentes à $\Sigma_i$ en $R_i$ se découpent suivant les deux hémisphères
$\in \boldsymbol{\varepsilon} = \pi_0(\Sigma \setminus \Sigma_{\mathbf{R}})$).
\marginal{NB Il faudrait introduire $\widetilde{\Sigma}_\varepsilon$ ($\varepsilon \in \boldsymbol{\varepsilon}$), les deux hémisphères sur $\widetilde{\Sigma}$, \uncertain{le choix} de $S_i^\varepsilon$ est \uncertain{caractérisé} par $S_i^\varepsilon \in \widetilde{\Sigma}_\varepsilon$.}

On a fait la figure dans le cas où
\[
  J = \{0, 1, \infty\}, \qquad \boldsymbol{\omega} = \{+1, -1\}
\]
($+$ donné par la perm. circulaire $1 \mapsto 2 \mapsto 3$, $-$ par l'inverse) et en
choisissant une orientation sur $E$
\[
  \boldsymbol{\varpi} \simeq \{+1, -1\}, \quad \text{d'où} \quad
  \boldsymbol{\varepsilon} \simeq \{+1, -1\} .
\]
\struck{On définit des chemins sur $\widetilde{\Sigma}$ \ill{} des parties de $\widetilde{\Sigma}$ (19) $a_i^{\omega}$, $b_i$, $c_i^{\omega}$ \ill{} (pour les $b_i^{\omega}$) \ill{} orientations \ill{} symétriques \ill{}}
\note{Bloc encadré et biffé au bas de la page ; il est repris, numérotation comprise, à la page suivante.}

\page{50}
\note{p. 5 de l'auteur.}
\[
  \text{\struck{(19)}}\qquad
  \boldsymbol{\varpi} \simeq \{\pm 1\}
\]
\struck{si besoin} i.e. on donne $\boldsymbol{\omega} \xrightarrow{\sim} \boldsymbol{\varepsilon}$,
donc donnant un procédé pour identifier l'une $\boldsymbol{\omega} = \Omega(J) = \Omega(\Sigma_{\mathbf{R}})$
\ill{} deux hémisphères.
\marginal{Il s'agit d'arcs \ill{} \uncertain{d'origine et d'extrémité} \ill{} \ill{} sur des cercles.}

\[
  (20)\quad
  \overset{12}{a_i^{\omega,\varepsilon}},\quad
  \overset{6}{b_i^{\omega}},\quad
  \overset{6}{c_i^{\varepsilon}},\quad
  \overset{6}{d_i^{\varepsilon}},
  \qquad i \in J,\ \omega \in \boldsymbol{\omega},\ \varepsilon \in \boldsymbol{\varepsilon}
\]
\note{Le numéro (20) est restitué : la page porte un numéro surchargé devant la ligne qui précède.}
On introduit sur $\widetilde{\Sigma}$ les \struck{\ill{}} chemins simples (segments orientés
plongés) ci-dessus ; $b, c, d$ chaque paquet forment un torseur sous $\mathfrak{S}_J$.
\marginal{NB $\operatorname{Rep}(J) \simeq J \times \boldsymbol{\omega}$ (21), $(i,j,k) \mapsto (i, \omega_{ijk})$.}
\marginal{Le paquet $(a_i^{\omega\varepsilon})$ \ill{} sous $\Gamma = \mathfrak{S}_J$ \ill{}.}
NB on a donc $g(a_i^{\omega}) = a_{g(i)}^{g(\omega)}$, etc. ; $\tau$ opérant par
\[
  (22)\quad
  \tau a_i^{\omega,\varepsilon} = a_i^{\omega,\varepsilon'}, \quad
  \tau b_i^{\omega} = b_i^{\omega}, \quad
  \tau c_i^{\varepsilon} = c_i^{\varepsilon'}, \quad
  \tau d_i^{\varepsilon} = d_i^{\varepsilon'} .
\]
On a
\[
  (23)\quad a_i^{\omega\varepsilon} : S_i^{\varepsilon} \longrightarrow R_i^{\omega}
  \quad \text{sur } U_i
\]
\struck{dans le sens de l'orientation de $\Sigma$ en $R_i$} (quart de tour), ce qui
caractérise $a_i^{\omega\varepsilon}$ \struck{explicitement}.
\[
  (24)\quad b_i^{\omega} : Q_i \longrightarrow R_{\omega' i}^{\omega'}
  \quad \text{sur } \widetilde{\Sigma}_{\mathbf{R}} \simeq \operatorname{Dic}(\Sigma_{\mathbf{R}}, J)
\]
\struck{(chemin d'origine \ill{})} (chemin sur $\widetilde{\Sigma}_{\mathbf{R}}$, d'origine
$Q_i$, dans la direction $\omega$, jusqu'à l'extrémité \struck{\ill{}} du segment
$\widetilde{\Sigma}_{\mathbf{R}} \simeq \operatorname{Dic}(\Sigma_{\mathbf{R}}, J)$, soit
$R_{\omega' i}^{\omega'}$),
\note{« Dic » tel quel sur la page.}
\[
  (25)\quad c_i^{\omega} : P^{\omega} \longrightarrow Q_i
  \qquad \text{suivant un arc de grand cercle}
\]
\[
  (26)\quad d_i^{\omega} : P^{\omega} \longrightarrow S_i^{\omega}
  \qquad \text{suivant un arc de grand cercle}
\]
\note{Dans (25) et (26) l'indice supérieur est écrit $\omega$ là où (20) porte $\varepsilon$.}
On pose aussi
\[
  (27)\quad b_i : R_{\omega i}^{\omega'} \longrightarrow R_{\omega' i}^{\omega'}, \qquad
  \widetilde{b}_i^{\omega} = b_i^{\omega} (b_i^{\omega'})^{-1}
\]
\marginal{chemin dans la direction $\omega$ \ill{} \struck{\ill{}}}
Tous les arcs $b_i^{\omega}$, $c_i^{\omega}$, $d_i^{\omega}$ se font sur des
\uncertain{quarts de} grands cercles de $\Sigma$, sauf que les extrémités des $c_i^{\omega}$,
$d_i^{\omega}$ ne sont pas dans $\Sigma^{*}$ mais dans
$\partial\widetilde{\Sigma} \setminus \Sigma^{*}$ ($U_i \subset \partial\widetilde{\Sigma}$).

\page{51}
\note{p. 6 de l'auteur (chiffre surchargé).}
Je vais introduire, dans la sphère $\widetilde{\Sigma}$ \uncertain{dont le degré de connexion est 0}, cinq ouverts, aussi grands que possible, qui soient
simplement connexes — obtenus en pratiquant des « incisions » qui relient entre eux les
\struck{trois} bords $U_i$ des trois « trous » \struck{\ill{}}. On pose
\[
  (27)\quad
  \left\{
  \begin{array}{ll}
    \widetilde{\Sigma}_\varepsilon = \widetilde{\Sigma} \setminus
      \bigcup_{i \in J} \operatorname{supp} d_i^{\varepsilon}
      & \varepsilon \in \boldsymbol{\varepsilon} \\[4pt]
    \widetilde{\Sigma}_i = \widetilde{\Sigma} \setminus
      \bigcup_{j \in J \setminus \{i\}} b_j &
  \end{array}
  \right.
\]
(où $b_j = \operatorname{supp} \widetilde{b}_j^{\omega}
= \operatorname{supp} c_j^{\omega} \cup \operatorname{supp} c_j^{\omega'}$ ).
\note{Le numéro (27) est employé deux fois de suite, p. 50 et ici.}
\drawing{deux sphères en marge. Sur la première, $\widetilde{\Sigma}_\varepsilon$ : $P_\varepsilon$ au sommet, les trois incisions $d$ menant aux points $S_i^{\varepsilon'}$, $S_j^{\varepsilon'}$, $S_k^{\varepsilon'}$ des trous, en trait épais. Sur la seconde, $\widetilde{\Sigma}_i$ : $Q_i$, $Q_j$, $R_i^{\omega}$ et les incisions le long de l'équateur.}

On voit bien que ces ouverts \struck{dans $\widetilde{\Sigma}$} sont simplement connexes,
connexes et \ill{} de la sphère ; dans $\widetilde{\Sigma}$ (car le complémentaire est
connexe). On dit que $P_\omega$ est le centre de $\widetilde{\Sigma}_\omega$, $Q_i$ le
centre de $\widetilde{\Sigma}_i$.
\marginal{en tant qu'ouverts \ill{}, et pour cette raison connexes \ill{}}

On pose
\[
  (28)\quad
  \left\{
  \begin{array}{l}
    \Pi_\varepsilon = \pi_1(\widetilde{\Sigma}, \widetilde{\Sigma}_\varepsilon)
      \simeq \pi_1(\Sigma^{*}, \Sigma^{*}_\varepsilon)
      \simeq \pi_1(\Sigma^{*}, P_\varepsilon) \\[4pt]
    \Pi_i = \pi_1(\widetilde{\Sigma}, \widetilde{\Sigma}_i)
      = \pi_1(\Sigma^{*}, \Sigma^{*}_i)
      \simeq \pi_1(\Sigma^{*}, Q_i)
  \end{array}
  \right.
\]
en faisant les identifications impliquées par ces formules. \struck{\ill{}} Pour tout
$x \in \widetilde{\Sigma}_\omega$ resp. $x \in \widetilde{\Sigma}_i$, on pose
\[
  (29)\quad
  \left\{
  \begin{array}{ll}
    \varphi_{\varepsilon,x} \text{ ou } \varphi_\omega :
      \Pi_\varepsilon \xrightarrow{\ \sim\ } \pi_1(\widetilde{\Sigma}, x)
      & \big(\simeq \pi_1(\Sigma^{*}, x) \text{ si } x \in \Sigma^{*}\big) \\[4pt]
    \varphi_{i,x} \text{ ou } \varphi_i :
      \Pi_i \xrightarrow{\ \sim\ } \pi_1(\widetilde{\Sigma}, x)
      & \text{—}
  \end{array}
  \right.
\]
\marginal{Les $\varphi$ \ill{} appelés \uncertain{isomorphismes de transition} \ill{}}

\page{52}
\note{Numéro de l'auteur « 6 », répété.}
l'isom. canonique déduit de l'inclusion. On a de plus des isomorphismes canoniques
\[
  (30)\quad
  \left\{
  \begin{array}{ll}
    \varphi_{i\varepsilon} : \Pi_\omega \longrightarrow \Pi_i
      & \text{défini par } c_i^{\varepsilon} : P_\varepsilon \to Q_i \\[4pt]
    \varphi_{\varepsilon i} : \Pi_i \longrightarrow \Pi_\omega
      & \text{---\ } (c_i^{\varepsilon})^{-1} : Q_i \to P_\varepsilon
  \end{array}
  \right.
\]
\marginal{NB on a $\varphi_{i\varepsilon} = \varphi_{Q_i,\varepsilon}$, $\varphi_{\varepsilon i} = \varphi_{P_\varepsilon, i}$.}
Par ces isom., les $\Pi_\varepsilon$, $\Pi_i$ sont reliés de façon transitive, mais bien
entendu pas par un syst. transitif d'isomorphismes — au contraire, les automorphismes de
l'un des groupes $\Pi_\varepsilon$, $\Pi_i$ obtenus en composant des isom. (30) suivant des
chemins fermés sur le graphe
\drawing{le graphe d'un octaèdre : $P_\varepsilon$ en haut, $P_{\varepsilon'}$ en bas, $Q_i$, $Q_j$, $Q_k$ sur l'équateur, chaque $P$ relié à chaque $Q$.}
sont tous les automorphismes intérieurs !

Les $\pi_1(\Sigma^{*}, \text{\ill{}})$ en les $P_\varepsilon$, $Q_i$ étant canonifiés comme
les $\Pi_\varepsilon$, $\Pi_i$, il reste à « repérer » les $\pi_1$ en les autres points
\uncertain{marqués} jusqu'à maintenant. Les choix naturels sont par
\struck{\ill{}}
\note{Une formule biffée suit, où l'on reconnaît $\Pi_\varepsilon$, $\Pi_i$ et des flèches ; elle est reprise en (31) à la page suivante.}

\page{53}
\[
  (31)\quad
  \left\{
  \begin{array}{ll}
    \varphi_{R_i^{\omega},\, \omega i} : \Pi_{\omega i} \xrightarrow{\ \sim\ }
      \pi_1(\widetilde{\Sigma}, R_i^{\omega})
      & \text{ou } \varphi_{\omega i} \\[4pt]
    \varphi_{S_i^{\varepsilon},\, \varepsilon} : \Pi_\varepsilon \xrightarrow{\ \sim\ }
      \pi_1(\widetilde{\Sigma}, S_i^{\varepsilon})
      & \text{ou } \varphi^{\varepsilon} \text{ ou } d_i^{\varepsilon}
  \end{array}
  \right.
\]
\note{Les indices de (31) sont lus avec doute.}
Le plus commode me semble de noter simplement $\varphi_i$, $\varphi_\varepsilon$, pour
préciser le groupe $\pi_1$ standard $\Pi_i$ ou $\Pi_\varepsilon$ qui est en jeu, le point
base $x$ ($= R_i^{\varepsilon}$ ou $S_i^{\varepsilon}$, disons) ne faisant pas de doute
dans le calcul.

\struck{Les arcs $a_i^{\omega}$ \ill{} sont trop longs, il faut les couper en deux en posant}
\[
  \text{\struck{$a_i^{\omega} = a_i^{\omega-} a_i^{\omega+}$, \quad $a_i^{\omega+} : R_i^{\omega} \to S_i^{\omega}$, \quad $a_i^{\omega-} : S_i^{\omega} \to R_i^{\omega'}$}}
\]
\note{Bloc encadré et biffé.}

Il faut enfin, pour des calculs \uncertain{pénibles}, donner une façon de repérer les
$\Pi_i$, $\Pi_\omega$ par un groupe standard \struck{$\Pi_{0,3}$}, pour lequel nous
prendrons
\[
  (32)\quad
  \Pi_{0,3} = \pi_1\big(\mathbf{P}^1_{\mathbf{C}} \setminus \{0,1,\infty\}
  \,[= U_{0,3}(\mathbf{C})],\ P^{+}\big)
  = \{\, l_0, l_1, l_\infty \mid l_\infty l_1 l_0 = 1 \,\}
\]
\[
  \subset \widetilde{\Pi}_{0,3}
  = \pi_1\big(U_{0,3}(\mathbf{C}),\ \mathfrak{S}_3 \times \{1,\tau\},\ P^{+}\big)
\]
\[
  = \{\, \rho, \sigma_\infty, \widetilde{\sigma}_\infty \mid
  \rho^3 = \sigma_\infty^2 = \widetilde{\sigma}_\infty^2 = 1,\
  \sigma_\infty(\rho) = \rho^{-1},\
  \widetilde{\sigma}_\infty(\sigma_\infty) = 1 \,\}
\]
\[
  \text{[i.e. } (\sigma_\infty \widetilde{\sigma}_\infty)^2 = 1 \text{]}
  \qquad \simeq \mathrm{GL}(2,\mathbf{Z})/\{\pm 1\}
\]
\struck{$\{\, \ldots \mid \rho^3 = \sigma_\infty^2 = 1 \,\} \simeq \mathrm{SL}(2,\mathbf{Z})/\pm 1$}
\note{La lecture de la relation $\widetilde{\sigma}_\infty(\sigma_\infty) = 1$ est douteuse ; le second membre de (32) est écrit au-dessus d'une première version biffée.}

\page{54}
\note{p. 7 de l'auteur.}
avec le formulaire familier dans $\widetilde{\Pi}_{0,3} \simeq \mathrm{SL}(2,\mathbf{Z})$ :
\[
  (33)\quad
  \left\{
  \begin{array}{l}
    \varepsilon_0 \overset{\mathrm{def}}{=} \sigma_\infty \rho, \quad
    \varepsilon_1 \overset{\mathrm{def}}{=} \rho(\varepsilon_0)
      \ (= \rho\, \varepsilon_0 \rho^{-1}), \\[4pt]
    \varepsilon_\infty \overset{\mathrm{def}}{=} \rho(\varepsilon_1)
      = \rho^2(\varepsilon_0) = \rho^{-1}(\varepsilon_0) \\[4pt]
    l_0 \overset{\mathrm{def}}{=} \varepsilon_0^2, \quad
    l_1 \overset{\mathrm{def}}{=} \varepsilon_1^2, \quad
    l_\infty \overset{\mathrm{def}}{=} \varepsilon_\infty^2, \quad
    l_\infty l_1 l_0 = 1 \\[4pt]
    \rho(l_i) = l_{\rho(i)}
      \ \ \{\rho(l_0) = l_1,\ \rho(l_1) = l_\infty,\ \rho(l_\infty) = l_0\}, \quad
    \rho(\varepsilon_i) = \varepsilon_{\rho(i)}
  \end{array}
  \right.
\]
NB on fait opérer $\rho$ sur $0, 1, \infty$ par $0 \mapsto 1 \mapsto \infty \mapsto 0$.
\note{Les $\varepsilon_i$ de (33) sont des éléments de $\widetilde{\Pi}_{0,3}$, sans rapport avec les $\varepsilon \in \boldsymbol{\varepsilon}$. La lecture $\varepsilon_0 = \sigma_\infty\rho$ est douteuse.}
\[
  (33)\quad
  \left\{
  \begin{array}{l}
    \tau_\infty \overset{\mathrm{def}}{=} \sigma_\infty \widetilde{\sigma}_\infty
      = \widetilde{\sigma}_\infty \sigma_\infty, \quad \tau_\infty^2 = 1 \\[4pt]
    \tau_\infty(l_0) = l_0^{-1}, \quad \tau_\infty(l_1) = l_1^{-1}, \quad
    \tau_\infty(l_\infty) = l_0^{-1} l_\infty^{-1} l_0 \\[4pt]
    \widetilde{\sigma}_\infty(l_0) = l_1^{-1}, \quad
    \widetilde{\sigma}_\infty(l_1) = l_0^{-1}, \quad
    \widetilde{\sigma}_\infty(l_\infty) = l_\infty^{-1} \\[4pt]
    \sigma_\infty(l_0) = l_1, \quad \sigma_\infty(l_1) = l_0, \quad
    \sigma_\infty(l_\infty) = l_1^{-1} l_0^{-1}
  \end{array}
  \right.
\]
\note{Le numéro (33) figure deux fois. La valeur de $\tau_\infty(l_\infty)$ est lue avec doute, et la dernière égalité porte une suite biffée illisible.}

Notons qu'un épinglage de la situation est fourni par un élément du produit
$\operatorname{Rep}(J, \boldsymbol{\varpi}) \simeq J \times \boldsymbol{\omega} \times \boldsymbol{\varpi}$,
donc par le choix d'un $(i, \omega, \varpi)$, définissant donc un
$\varepsilon = \omega \wedge \varpi \in \boldsymbol{\varepsilon} \simeq \boldsymbol{\omega} \wedge \boldsymbol{\varpi}$.
On va définir pour $r = (i, \omega, \varpi)$ des isom.
\[
  (34)\quad
  \left\{
  \begin{array}{l}
    \psi_\varepsilon^{i,\omega} : \Pi_{0,3} \xrightarrow{\ \sim\ } \Pi_\varepsilon \\[4pt]
    \psi_i^{\omega,\varepsilon} : \Pi_{0,3} \xrightarrow{\ \sim\ } \Pi_i
  \end{array}
  \right.
\]
\marginal{isom. dits d'épinglage}
donnant lieu comme de juste à un triangle commutatif (35)

\begin{tikzcd}[column sep=large]
  & \Pi_{0,3} \arrow[dl, "\psi_\varepsilon^{i\omega}"'] \arrow[dr, "\psi_i^{\omega,\varepsilon}"] & \\
  \Pi_\varepsilon \arrow[rr, bend left=15, "\varphi_{i\varepsilon}"] & & \Pi_i \arrow[ll, bend left=15, "\varphi_{\varepsilon i}"]
\end{tikzcd}

\page{55}
ce qui fait qu'il suffit de définir l'un des deux isom. (34) — et vu la définition de
$\Pi_{0,3}$ dans (32), il est naturel de définir plutôt l'iso
\[
  \text{\struck{(35)}}\quad
  \psi_r = \psi_{\varepsilon}^{i,\omega} : \Pi_{0,3} \xrightarrow{\ \sim\ } \Pi_\varepsilon
  \qquad r = (i, \omega, \varpi)
\]
par transport de structure comme
\[
  (36)\quad
  \psi_\varepsilon^{i,\omega} = \pi_1(u_{i,\omega,\varepsilon} ; P^{+}),
  \quad \text{où } u_r = u_{i,\omega,\varepsilon} = u_\varepsilon^{i,\omega} :
\]
\[
  \big(U_{0,3}(\mathbf{C}), P^{+}\big) \longrightarrow (\Sigma^{*}, P_\omega)
\]
est l'unique iso. conforme qui transforme le repère standard
$\{\infty, \omega_{01\infty}, +1\}$ \struck{\ill{}} dans le repère \struck{\ill{}}
$(i, \omega, \varepsilon)$.
\note{Le numéro (36) est restitué : la page porte un numéro surchargé, lu « (35) » ou « (36) ».}
\marginal{Au couple $(i,\omega)$ on associe le repère $r = (\omega i, \omega' i, i)$ de $J$ \struck{\ill{}}, i.e. $0 \mapsto \omega i$, $1 \mapsto \omega' i$, $\infty \mapsto i$. \struck{$r = \ldots$} \uncertain{Avec ces notations} $\rho_\omega(r) = r \circ \rho = (\omega i, \omega)$.}

Pour un $P_\varepsilon$ fixé, il y a six choix pour les
$\psi_r = \psi_{\varepsilon}^{i,\omega} : \Pi_{0,3} \to \Pi_\varepsilon$ — et pour un $Q_i$
fixé, il y a quatre choix pour $\omega, \varepsilon$. Cela fait \add{trois resp.}
\struck{douze} choix si on suppose qu'on a choisi un élément de $\boldsymbol{\varpi}$, i.e.
$\Sigma$ est orienté, d'où $\boldsymbol{\varepsilon} \simeq \boldsymbol{\omega}$, et si on se
borne aux $\varepsilon$ qui correspondent aux $\omega$) — il faudrait donc voir comment
passer \add{d'un choix} de l'un à l'autre.

\page{56}
\note{Numéro de l'auteur « 9 » au-dessus de la première ligne, lu avec doute ; la p. 8 manque à la suite.}
Notons que l'isom.
\[
  \psi_r : \Pi_{0,3} \xrightarrow{\ \sim\ } \Pi_\omega, \qquad
  r \in \operatorname{Rep}(J, \boldsymbol{\varpi})
  \simeq \operatorname{Isom}\big((\Sigma_0^{*}, P_{+}), (\Sigma^{*}, P_{+})\big)
\]
est compatible avec le transport de structure, \struck{en ce qui} \add{dans le sens} qu'il
élément
\[
  g \in \widetilde{\mathfrak{S}}_J \subset \Gamma = \check{\mathfrak{S}}_J
\]
\[
  \Gamma_{P_\varepsilon}
  = \{\, (g, \tau^{\alpha}) \in \mathfrak{S}_J \times \mathbf{Z}/2\mathbf{Z} \mid
  \alpha = 0 \text{ si } \operatorname{sg} g = 1,\
  \alpha = 1 \text{ si } \operatorname{sg} g = -1 \,\}
\]
\[
  = \{\, g \in \Gamma \mid g(\varepsilon) = \varepsilon \,\}
\]
\note{Le signe « $=$ » vertical relie $\widetilde{\mathfrak{S}}_J$ à $\Gamma_{P_\varepsilon}$ ; l'accent sur le second $\mathfrak{S}_J$ est lu avec doute.}
on aura commutativité dans

\begin{tikzcd}[column sep=small]
  & \Pi_{0,3} \arrow[dl, "{\psi_r = \psi_{\omega,i}^{\varepsilon}}"'] \arrow[dr, "{\psi_{g(r)} = \psi_{g(\omega),g(i)}^{g(\varepsilon)}}"] & \\
  \Pi_\varepsilon \arrow[rr, "\pi_1(g)"'] & & \Pi_\varepsilon
\end{tikzcd}

\[
  (37)\qquad
  \pi_1(g) \text{ par transport de structure}
\]

ce qui suffit au passage d'un $\psi_r$ à $\psi_{g(r)} = \pi_1(g) \circ \psi_r$, vu que
$\widetilde{\mathfrak{S}}_J$ est simplement transitif sur l'ensemble des repères $r$
correspondant à un $\varepsilon$ fixé. Mais on aimerait exprimer le passage par une
formule de transformation \emph{dans} $\Pi_{0,3}$. Notons que la donnée des repères
définit un isom. canonique (encore noté $u_r$)

\page{57}
\[
  u_r : \widetilde{\mathfrak{S}}_3 \xrightarrow{\ \sim\ } \widetilde{\mathfrak{S}}_3
\]
et tout repère $r'$ \add{\ill{} $= \varepsilon$ pour $r$,} s'écrit de façon unique
\[
  r' = r g_0 \quad \text{avec } g_0 \in \widetilde{\mathfrak{S}}_3
\]
\[
  (38)\quad \text{i.e. } u_{r'} = u_r u_{g_0}
  \qquad (\text{avec } u_r(g_0) = g \text{ mais peu importe})
\]
d'où, en passant à l'action sur les $\pi_1$ en $P_{+}$ resp. $P_\varepsilon$, on trouve
\[
  (39)\quad \psi_{r'} = \psi_r \circ \pi_1(g_0, P_{+})
  \qquad \text{si } r' = r \circ g_0
\]
\marginal{$\psi_r$, $\psi_{r'}$ désignent les isom. d'épinglage \ill{} $P_\varepsilon$ \ill{} $\psi_\varepsilon(l_0)$, $\psi_\varepsilon(l_1)$ \ill{} les lacets standard \ill{} $Q_i$ \ill{} $R_{\omega i}$.}
\marginal{NB \ill{} les automorphismes \ill{} $P_\varepsilon$, $R_{\omega i}$ \ill{} d'orientation $\boldsymbol{\varpi} = \omega \wedge \varepsilon$.}
\drawing{en marge, un petit schéma : $P_\varepsilon$ relié par un lacet $l_0$ au trou $R_{\omega i} \simeq R_0$ (d'orientation $\varpi = \omega \wedge \varepsilon$), au point $Q_\infty$ et au trou $R_{\omega' i} \simeq R_1$.}

Explicitons ceci pour
\[
  g_0 = \rho \quad \text{donc} \quad
  r' = r\rho = (\omega i, \omega, \varpi) \quad \text{si } r = (i, \omega, \varpi)
\]
on trouve
\[
  (40)\quad
  \psi_\varepsilon^{\omega i, \omega}
  = \overbrace{\psi_\varepsilon^{i,\omega} \circ \rho}^{\rho_\omega \circ \psi_\varepsilon^{i,\omega} =}
\]
où on a identifié $\rho$ aussi à un élément de $\widetilde{\Pi}_{0,3}$,
\marginal{NB $\rho_\omega$ est la rotation autour de l'axe de $P_\varepsilon$, $P_{\varepsilon'}$, d'un tiers de tour, \struck{\ill{} d'orientation} \ill{} défini par $\omega$ sur l'\uncertain{équateur} \ill{} dans $\Pi_{0,3}$.}
On trouve de même, pour un \uncertain{automorphisme} qui renverse l'orientation de
$\Sigma$, i.e. pour un
\[
  \widetilde{\sigma}_\infty : \quad
  0 \mapsto 1, \quad 1 \mapsto 0, \quad \infty \mapsto \infty,
  \qquad \varpi \mapsto \varpi'
\]
\[
  g_0 = \widetilde{\sigma}_\infty, \qquad
  r' = r \sigma_\infty = (\omega' i, \omega', \varpi')
\]
d'où
\[
  (41)\quad
  \psi_\varepsilon^{\omega' i, \omega'}
  = \overbrace{\psi_\varepsilon^{i,\omega} \circ \widetilde{\sigma}_\infty}^{\widetilde{\sigma}_{?} \circ \psi_\varepsilon^{i,\omega} =}
\]
($\widetilde{\sigma}_i$ est la symétrie p.r. aux grands cercles méridiens passant par
$Q_i$).
\note{Dans l'accolade de (41) l'indice de $\widetilde{\sigma}$ est illisible ; le premier terme de $r'$ est surchargé.}

\page{58}
\note{p. 10 de l'auteur.}
Passons au cas de $Q_i$ — la définition \struck{\ill{}} des
$\psi_r = \psi_i^{\omega,\varepsilon} : \Pi_{0,3} \to \Pi_i$ \struck{\ill{}} est
\[
  \psi_r = \psi_i^{\omega,\varepsilon}
  = \pi_1(\underbrace{u_{i,\omega,\varepsilon}}_{u_r} ; Q_\infty)
\]
où de même $u_r = u_{i,\omega,\varepsilon}$ est l'unique isom. de $U_{0,3}$
\add{dans $\Sigma^{*}$} qui transforme le repère \struck{fixé} standard
\[
  (\infty,\ \omega_0 = +1,\ \varpi_0 = +1 \text{ donc } \varepsilon = \omega_0 \wedge \varpi_0 = +1)
\]
en $r$, i.e. qui transforme $P_{+}$ en $P_\varepsilon$, $Q_0$ en $Q_i$, $Q_1$ en
$Q_{\omega i}$. Les épinglages $\psi_i^{\omega,\varepsilon}$ ($i$ fixé, $\omega, \varepsilon$
variables) \struck{\ill{}} de $\Pi_i$ correspondant aux $u_r$ qui transforment $Q_0$ en
$Q_i$, \uncertain{se sont déterminés} par compatibilité avec des éléments
$g_0 \in \Gamma_0 = \mathfrak{S}_3 \times \{1, \tau\}$ qui stabilisent $Q_0$ — ces éléments
\struck{\ill{}} forment un groupe $\simeq \mathbf{Z}/2 \times \mathbf{Z}/2$ engendré par
$\sigma_\infty, \tau$. On trouve encore
\[
  \psi_{r'} = \psi_r \circ \pi_1(g_0 ; Q_\infty) \qquad \text{si } r' = r \circ g_0
\]
où ici $\psi_r$, $\psi_{r'}$ désignent les iso. d'épinglage relatifs à $Q_i$. Il suffit
d'expliciter ces formules dans les cas de $g_0 = \sigma_\infty$, $g_0 = \tau$, ce qui
revient à calculer leur effet sur \struck{\ill{}} $\pi_1(U_{0,3} ; Q_\infty)$, quand on
identifie celui-ci à $\Pi_{0,3} = \pi_1(U_{0,3} ; P_{+})$, grâce au chemin $d_\infty^{+}$ de
$P_{+}$ à $Q_\infty$. On sait que ces générateurs
\marginal{NB on identifie $\pi_1(U_{0,3} ; Q_\infty)$ à $\Pi_{0,3} \overset{\mathrm{def}}{=} \pi_1(U_{0,3} ; P_{+})$ à l'aide du chemin $d_\infty^{+}$ de $P_{+}$ à $Q_\infty$.}
\note{Les $Q_i$ de cette page jouent le rôle que (31) donnait aux $\Pi_i$ ; le chemin $d_\infty^{+}$ joint ici $P_{+}$ à $Q_\infty$, alors qu'en (25)–(26) c'est $c$ qui mène aux $Q_i$ et $d$ aux $S_i$.}

\page{59}
sont $\operatorname{int}(\sigma_\infty)$, $\operatorname{int}(\tau_\infty)$. Or dans le cas
$g_0 = \tau_\infty$ resp. $g_0 = \sigma_\infty$, \struck{donc} on trouve
$r' = (i, \omega, \varepsilon')$ resp. $r' = (i, \omega', \varepsilon')$, donc
\[
  (39)\quad
  \left\{
  \begin{array}{l}
    \psi_i^{\omega,\varepsilon'} = \tau \circ \psi_i^{\omega,\varepsilon}
      = \psi_i^{\omega,\varepsilon} \circ \tau_\infty \\[4pt]
    \psi_i^{\omega',\varepsilon'} = \sigma_i \circ \psi_i^{\omega,\varepsilon}
      = \psi_i^{\omega,\varepsilon} \circ \sigma_\infty
  \end{array}
  \right.
\]
($\sigma_i$ : transposition $\in \mathfrak{S}_J$ qui fixe $i$) dont en composant
\[
  (40)\quad
  \psi_i^{\omega',\varepsilon}
  = \widetilde{\sigma}_i \circ \psi_i^{\omega,\varepsilon}
  = \psi_i^{\omega,\varepsilon} \circ \widetilde{\sigma}_\infty
\]
\note{Sous $\widetilde{\sigma}_i$ la page porte « $\tau\sigma_i$, symétrie autour des grands cercles passant par $P_\varepsilon$, $P_{\varepsilon'}$, $Q_i$ ». Les numéros (39) et (40) sont ici repris une seconde fois, après (39)–(41) de la p. 57.}

Dans tous ces calculs, suivant sans doute d'anciennes habitudes, j'ai privilégié les
$P_\varepsilon$ et les $Q_i$ comme pts base — alors que pour le travail que j'ai en vue,
ce sont les points $R_i^{\omega}$ qui sont les plus cruciaux, où les calculs \ill{} en ces
pts expriment directement les phénomènes de « monodromie » autour des « trous » $R_i$.
Je vais réparer cette négligence, en introduisant, en

\page{60}
\note{p. 11 de l'auteur.}
plus des $\widetilde{\Sigma}_\omega$, $\Pi_\omega$ et $\widetilde{\Sigma}_i$, $\Pi_i$, des
régions
\[
  \widetilde{\Sigma}_r = \widetilde{\Sigma}_i^{\omega,\varepsilon}
\]
associées aux repères $r = (i, \omega, \struck{\ill{}}\, \varepsilon)$ de
$(J, \boldsymbol{\varepsilon})$ (ou ce qui revient au même, de $(J, \boldsymbol{\varpi})$),
qu'on peut identifier, géométriquement, aux \add{douze} triangles-drapeaux
$D_r = P_\varepsilon Q_i R_{\omega i}$, correspondant à la subdivision barycentrique
standard de la sphère $\Sigma$ définie par l'équateur $\Sigma_{\mathbf{R}}$ et par les trois
points $R_i$ sur celui-ci. Le « \uncertain{trou} » \add{sur $\widetilde{\Sigma}$} \ill{}
issu de l'un de ces triangles est un quadrilatère $\widetilde{D}_r$
\[
  P_\varepsilon Q_{\omega' i} R_i^{\omega} S_i^{\varepsilon}
  \simeq P_{+} Q_\infty R_0^{+} S_0^{+} .
\]
\drawing{une sphère : $P_\varepsilon \simeq P_{+}$ au sommet, $P_{\varepsilon'}$ en bas ; le quadrilatère $\widetilde{D}_r$ hachuré entre le méridien de $Q_i \simeq Q_\infty$ et celui du trou $R_{\omega i} \simeq R_1$ ; autour des trous les points $R_{\omega i}^{\omega}$, $S_i^{\varepsilon} \simeq S_0^{\varepsilon}$, $Q_{\omega' i}$ ; légende $r = (i, \omega, \varepsilon)$, $J = \ldots$}

On aura à travailler, au « voisinage immédiat » de ce repère \struck{\ill{}}, au
\struck{triangle-drapeau} \add{drapeau} $\widetilde{\Sigma}_r$, avec les symétriques de
$\Sigma_r$ p.r. à ses trois côtés, et ses transformés par $\Sigma_r$,
\struck{ce qui conduit} à regarder la réunion de $\Sigma_r$ avec ses trois transformés
\[
  \tau(\Sigma_r), \quad \widetilde{\sigma}_i(\Sigma_r), \quad
  \widetilde{\sigma}_{\omega i}(\Sigma_r)
\]
(qui correspondent, sur la sphère standard avec le repère correspondant au triangle
$P_{+} Q_\infty R_0$, à $\tau$, $\widetilde{\sigma}_\infty$, $\widetilde{\sigma}_0$) et sur
$\widetilde{\Sigma}$ à prendre
\[
  (41)\quad
  \widetilde{\Sigma}_r = \widetilde{\Sigma}_i^{\omega,\varepsilon}
  = \widetilde{D}_r \cup \tau(\widetilde{D}_r) \cup \widetilde{\sigma}_i(\widetilde{D}_r)
  \cup \widetilde{\sigma}_{\omega i}(\Sigma_r)
\]
limité par le contour
\[
  P_\varepsilon Q_{\omega' i} R_{\omega' i}^{\omega} S_i^{\varepsilon} \ldots
  P_{\varepsilon'} \ldots Q_i R_{\omega' i}^{\omega'} S_{\omega' i}^{\varepsilon} P_\varepsilon
\]
\note{Le contour est lu en partie seulement : plusieurs sommets intermédiaires, serrés en bas de page, restent illisibles et sont remplacés par « $\ldots$ ». L'argument se poursuit au-delà de ce lot.}

\end{document}
